\documentclass[a4paper,fontset=windows]{ctexart}

\usepackage{anyfontsize}
\usepackage{amsmath,amssymb,mathrsfs}
\usepackage{autobreak}
\allowdisplaybreaks
\usepackage{indentfirst}
\setlength{\parindent}{0em}
\usepackage{geometry}
\geometry{top=3cm,bottom=3cm,left=2.75cm,right=2.75cm}
\usepackage{setspace}
\usepackage{bm}
\usepackage{graphicx,caption,subcaption}
\usepackage{indentfirst}
\usepackage{multirow,longtable,booktabs}
\renewcommand{\arraystretch}{1.25}


\begin{document}

\title{\textbf{实验三十\ \ 用示波器观测动态磁滞回线}}
\author{1900011604 张植竣}
\date{2021年6月14日}

\maketitle

\section*{1\ \ 实验数据和现象}

\subsection*{1.1\ \ 室温下的测量数据}

\begin{table}[!ht]
    \begin{center}
        \begin{tabular}{|c|c|c|c|c|c|}
            \hline
            \multicolumn{2}{|c|}{铂电阻温度计} & \multicolumn{2}{|c|}{硅二极管温度计} & \multicolumn{2}{|c|}{超导样品} \\
            \hline
            $U_\mathrm{Pt}/\mathrm{mV}$ & $I_\mathrm{Pt}/\mathrm{mA}$ & $U_\mathrm{Si}/\mathrm{V}$ & $I_\mathrm{Si}/\mathrm{mA}$ & $U_\mathrm{YBCO}/\mathrm{mV}$ & $I_\mathrm{YBCO}/\mathrm{mA}$ \\
            \hline
            110.64 & 1.0000 & 0.5034 & 0.10003 & 0.078 & 10.0159 \\
            \hline
        \end{tabular}
    \end{center}
\end{table}

室温：$300.29\,\mathrm{K}=27.14\text{\textcelsius}$

超导样品电阻：$7.79\,\mathrm{m}\Omega$

\subsection*{1.2\ \ 低温温度计的比对}

{\small
\begin{center}
    \begin{longtable}{cccccc}
        \toprule[0.1em]
        $U_\mathrm{Pt}/\mathrm{mV}$ & $U_\mathrm{Si}/\mathrm{V}$ & $U_\mathrm{TC}/\mathrm{mV}$ & $T/\mathrm{K}$ & $R_\mathrm{Pt}/\Omega$ & $R_\mathrm{Si}/\mathrm{k}\Omega$ \\
        \midrule[0.06em]
        88.96 & 0.6395 & 3.901 & 244.93 & 88.96 & 6.393 \\
        87.24 & 0.6509 & 3.739 & 240.59 & 87.24 & 6.507 \\
        85.84 & 0.6601 & 3.621 & 237.05 & 85.84 & 6.599 \\
        84.25 & 0.6707 & 3.511 & 233.04 & 84.25 & 6.705 \\
        82.97 & 0.6792 & 3.384 & 229.81 & 82.97 & 6.790 \\
        81.49 & 0.6893 & 3.258 & 226.08 & 81.49 & 6.891 \\
        80.00 & 0.6994 & 3.156 & 222.34 & 80.00 & 6.992 \\
        78.46 & 0.7101 & 3.309 & 218.48 & 78.46 & 7.099 \\
        76.98 & 0.7202 & 3.193 & 214.77 & 76.98 & 7.200 \\
        74.79 & 0.7351 & 3.026 & 209.29 & 74.79 & 7.349 \\
        72.96 & 0.7477 & 2.889 & 204.73 & 72.96 & 7.475 \\
        71.10 & 0.7604 & 2.752 & 200.09 & 71.10 & 7.602 \\
        69.67 & 0.7701 & 2.648 & 196.53 & 69.67 & 7.699 \\
        68.39 & 0.7789 & 2.552 & 193.35 & 68.39 & 7.787 \\
        66.97 & 0.7885 & 2.452 & 189.83 & 66.97 & 7.883 \\
        65.61 & 0.7977 & 2.357 & 186.45 & 65.61 & 7.975 \\
        64.89 & 0.8026 & 2.307 & 184.67 & 64.89 & 8.024 \\
        62.95 & 0.8157 & 2.175 & 179.87 & 62.95 & 8.155 \\
        61.48 & 0.8255 & 2.077 & 176.24 & 61.48 & 8.253 \\
        60.01 & 0.8352 & 1.981 & 172.61 & 60.01 & 8.349 \\
        58.51 & 0.8451 & 1.884 & 168.92 & 58.51 & 8.448 \\
        56.89 & 0.8556 & 1.782 & 164.94 & 56.89 & 8.553 \\
        55.46 & 0.8649 & 1.693 & 161.43 & 55.46 & 8.646 \\
        53.62 & 0.8767 & 1.582 & 156.93 & 53.62 & 8.764 \\
        52.47 & 0.8840 & 1.514 & 154.12 & 52.47 & 8.837 \\
        51.00 & 0.8933 & 1.427 & 150.53 & 51.00 & 8.930 \\
        49.51 & 0.9026 & 1.341 & 146.90 & 49.51 & 9.023 \\
        49.01 & 0.9057 & 1.314 & 145.68 & 49.01 & 9.054 \\
        48.11 & 0.9113 & 1.262 & 143.51 & 48.11 & 9.110 \\
        47.00 & 0.9182 & 1.197 & 140.81 & 47.00 & 9.179 \\
        45.98 & 0.9245 & 1.141 & 138.34 & 45.98 & 9.242 \\
        44.93 & 0.9309 & 1.084 & 135.80 & 44.93 & 9.306 \\
        43.99 & 0.9366 & 1.035 & 133.54 & 43.99 & 9.363 \\
        \bottomrule[0.1em]
    \end{longtable}
\end{center}
}

硅二极管温度计电阻随温度的变化规律：
\[
    y = 2.059\times10^{-7}x^3 - 0.000124x^2 - 0.002456x + 11.41,\quad
    R^2 = 0.9999979974
\]

\begin{figure}[!htbp]
    \centering
    \includegraphics[width=0.75\linewidth]{1.png}
    \caption{硅二极管温度计电阻随温度的变化规律}
\end{figure}

温差电偶温度计电压随温度的变化规律：
\[
    y = 4.559\times10^{-9}x^3 - 5.343\times10^{-5}x^2 - 0.007546x - 0.936,\quad
    R^2 = 0.9999957704
\]

\begin{figure}[!htbp]
    \centering
    \includegraphics[width=0.75\linewidth]{2.png}
    \caption{温差电偶温度计电压随温度的变化规律}
\end{figure}

\subsection*{1.3\ \ 超导转变曲线的测量}

{\small
\begin{center}
    \begin{longtable}{ccccccccc}
        \toprule[0.1em]
        $U_\mathrm{Pt}/\mathrm{mV}$ & $U_\mathrm{YBCO}/\mathrm{mV}$ & $T/\mathrm{K}$ & $U_\mathrm{Pt}/\mathrm{mV}$ & $U_\mathrm{YBCO}/\mathrm{mV}$ & $T/\mathrm{K}$ & $U_\mathrm{Pt}/\mathrm{mV}$ & $U_\mathrm{YBCO}/\mathrm{mV}$ & $T/\mathrm{K}$ \\
        \midrule[0.06em]
        43.01 & 0.040 & 131.17 & 28.63 & 0.029 & 96.96 & 26.44 & 0.014 & 91.82 \\
        42.00 & 0.040 & 128.74 & 28.01 & 0.028 & 95.50 & 26.43 & 0.013 & 91.80 \\
        41.00 & 0.039 & 126.34 & 27.57 & 0.027 & 94.47 & 26.43 & 0.012 & 91.80 \\
        40.01 & 0.038 & 123.96 & 27.17 & 0.026 & 93.53 & 26.42 & 0.011 & 91.77 \\
        39.00 & 0.038 & 121.55 & 26.94 & 0.025 & 92.99 & 26.41 & 0.010 & 91.75 \\
        38.00 & 0.037 & 119.16 & 26.75 & 0.024 & 92.55 & 26.40 & 0.009 & 91.73 \\
        37.00 & 0.037 & 116.77 & 26.63 & 0.023 & 92.27 & 26.40 & 0.008 & 91.73 \\
        35.99 & 0.037 & 114.37 & 26.58 & 0.022 & 92.15 & 26.40 & 0.007 & 91.73 \\
        34.96 & 0.036 & 111.91 & 26.54 & 0.021 & 92.05 & 26.39 & 0.006 & 91.70 \\
        34.00 & 0.035 & 109.63 & 26.52 & 0.020 & 92.01 & 26.38 & 0.005 & 91.68 \\
        32.97 & 0.034 & 107.20 & 26.49 & 0.019 & 91.94 & 26.37 & 0.004 & 91.66 \\
        31.96 & 0.033 & 104.81 & 26.49 & 0.018 & 91.94 & 26.36 & 0.003 & 91.63 \\
        30.96 & 0.032 & 102.45 & 26.47 & 0.017 & 91.89 & 26.32 & 0.002 & 91.54 \\
        29.96 & 0.031 & 100.09 & 26.46 & 0.016 & 91.87 & 26.28 & 0.001 & 91.45 \\
        29.16 & 0.030 & 098.21 & 26.45 & 0.015 & 91.84 & 26.23 & 0.000 & 91.33 \\
        \bottomrule[0.1em]
    \end{longtable}
\end{center}
}

\begin{figure}[!htbp]
    \centering
    \includegraphics[width=0.75\linewidth]{3.png}
    \caption{最小二乘法处理超导样品测量数据}
\end{figure}

最小二乘法处理起始转变温度之前一段的温度变化规律$R_\mathrm{n}T$：
\[
    R_\mathrm{n} = 0.04196T - 1.1046,\quad r = 0.9999807754
\]

\begin{figure}[!htbp]
    \centering
    \includegraphics[width=0.75\linewidth]{4.png}
    \caption{超导样品电阻随温度的变化规律}
\end{figure}

超导样品的转变温度：$91.81\,\mathrm{K}$

\subsection*{1.4\ \ 液氮沸点下的测量数据}

\begin{table}[!ht]
    \begin{center}
        \begin{tabular}{|c|c|c|c|c|c|}
            \hline
            \multicolumn{2}{|c|}{铂电阻温度计} & \multicolumn{2}{|c|}{硅二极管温度计} & \multicolumn{2}{|c|}{超导样品} \\
            \hline
            $U'_\mathrm{Pt}/\mathrm{mV}$ & $I'_\mathrm{Pt}/\mathrm{mA}$ & $U'_\mathrm{Si}/\mathrm{V}$ & $I'_\mathrm{Si}/\mathrm{mA}$ & $U'_\mathrm{YBCO}/\mathrm{mV}$ & $I'_\mathrm{YBCO}/\mathrm{mA}$ \\
            \hline
            20.32 & 0.9998 & 1.0702 & 0.10017 & 0.000 & 10.0177 \\
            \hline
        \end{tabular}
    \end{center}
\end{table}

液氮沸点温度：$77.49\,\mathrm{K}$

超导样品的电阻：$0.00\,\mathrm{m}\Omega$

乱真电动势：$0.000\,\mathrm{mV}$

铂电阻温度计电流的相对改变：
\[
    \frac{\Delta I_\mathrm{Pt}}{I_\mathrm{Pt}} = \frac{\left|I'_\mathrm{Pt} - I'_\mathrm{Pt}\right|}{I_\mathrm{Pt}} = 2.0\times10^{-4}
\]
硅二极管温度计电流的相对改变：
\[
    \frac{\Delta I_\mathrm{Si}}{I_\mathrm{Si}} = \frac{\left|I'_\mathrm{Si} - I'_\mathrm{Si}\right|}{I_\mathrm{Si}} = 1.4\times10^{-3}
\]
超导样品电流的相对改变：
\[
    \frac{\Delta I_\mathrm{YBCO}}{I_\mathrm{YBCO}} = \frac{\left|I'_\mathrm{YBCO} - I'_\mathrm{YBCO}\right|}{I_\mathrm{YBCO}} = 1.8\times10^{-4}
\]
可见铂电阻温度计和超导样品的电路稳定性约在万分之一量级，硅二极管温度计的电路稳定性在千分之一量级。

\end{document}